\clearpage
\fancyhead[L]{\small OBSERVATIONS ASTRONOMIQUES.}
\hypertarget{native-646-lower}{}
\begin{center}\small Göttingische gelehrte Anzeigen. No. 129. pp. 1281..1284. 14 August 1826.\\[3pt]\rule{.72\linewidth}{.3pt}\end{center}

\emph{Observations astronomiques faites à l'observatoire royal de Paris, publiées par le bureau des longitudes.} Paris, 1825. From Bachelier. 402 pages in folio.

The observations made at the Paris Observatory in the years 1800..1809 were published in the volumes of the Connaissance des tems for 1808..1812 and 1823..1825. By a decision of the Bureau des longitudes, the later observations are to be printed separately, and this first volume contains, with the exception of the observations made with \textsc{Reichenbach}'s repeating circle of three feet in diameter, those made in the years 1810..1819 by four observers, Messrs. \textsc{Bouvard}, \textsc{Arago}, \textsc{Mathieu} and \textsc{Nicollet}. The instruments used for the observations are described in the introduction. The transit telescope, begun by \textsc{Ramsden} and completed by \textsc{Berge}, was set up in August 1803. It has a focal length of $7\frac{1}{2}$ feet and an aperture of 4 inches. The axis is 4 feet long, and the magnification is not quite a hundredfold. Lateral movement of the eyepiece and illumination through the axis, as is now generally customary. To adjust the transit telescope, two meridian marks have been set up, the northern one on the Palais du Luxembourg at a distance of 1364 metres and the southern one on a pyramid in the plain of Montrouge, 1840 metres away. The targets are circular holes $2\frac{1}{4}$ inches in diameter in metal plates that can be shifted horizontally. The northern hole appears against a white-painted sheet of iron placed behind it, the southern one against the sky. The apparent diameters are $9''07$ and $7''1$. There are five vertical threads. Their intervals are taken to be exactly equal, and stars on the equator cross them in $17^{\rm s}36$. The material of the threads is not mentioned. The observations with this instrument take up more than half the volume. They concern the Sun, the Moon, most of the planets and \textsc{Maskelyne}'s fixed stars, with other objects rarely observed. Some lower culminations of the comets of 1811 and 1819 are recorded. As for the observations themselves, they fall somewhat short of what modern astronomy requires of independent observations with instruments of such remarkable dimensions. They afford no adequate means of determining how accurately the meridian marks lie in the meridian plane. Given the short distance of these marks and the long focal length of the telescope, the marks must show considerable parallax relative to the threads and appear poorly defined. Moreover, they were generally used for checks only infrequently. In 1810, for example, the first such use was on 11 May, after it had often been noted in the earlier months that fog had kept them from being seen. If the unfavourable site allowed the meridian marks to be used so rarely, frequent observations of circumpolar stars would have been doubly necessary. Yet even the pole star was seldom observed, and successive culminations \SourceLeaf{647}above and below the pole still less often. The axis was indeed often checked to see whether it was horizontal. But there was no means at all of checking whether its two pivots had the same diameter, which one must not assume even in instruments made by the foremost instrument makers. The perpendicularity of the optical axis to the axis of rotation was also checked only rarely. A great many observations of stellar occultations are interspersed in this log of observations with the transit telescope. It is all the more regrettable that the circumstances described always leave some slight uncertainty in the determination of absolute time, since astronomers are accustomed always to refer their statements of geographical longitude to the Paris Observatory.

\begingroup\rightskip=0pt plus 1em The second section of the work consists of observations made with \textsc{Bird}'s $7\frac{1}{2}$-foot mural quadrant. The telescope is likewise $7\frac{1}{2}$ feet long, has an aperture of $2\frac{1}{4}$ inches and magnifies 70..80 times. The objects observed are the same as with the transit telescope. An attempt was made to determine the collimation error at a number of points from numerous observations of fixed stars, taking as the basis for their declinations a mean of the values given by several astronomers. There is no need to say what standing the observations made with this instrument have in the present state of practical astronomy. This quadrant serves for the southern half of the meridian. On the north side hangs a second, five-foot quadrant by \textsc{Sisson}, the very one that \textsc{Lalande} once used in Berlin. This volume, however, contains no observations made with that instrument. --- A mural circle by \textsc{Fortin} has been in use since 1823. We may expect its description in the next volume.\par\endgroup

The third section contains observations made with \textsc{Bellet}'s parallactic instrument, which is set up in a separate part of the observatory. Each of its two circles is 13 inches in diameter. The telescope is three feet long, has a focal length of nearly $2\frac{1}{2}$ inches and magnifies 40..50 times. The observations concern in part the comets that appeared in the period 1810..1819 and in part the lunar spot Manilius, observed to investigate the Moon's libration. As is well known, the very valuable results of this investigation were published several years ago in the Connaissance des tems. This instrument was transferred to the observatory in Marseille in 1823 and replaced by a large equatorial by \textsc{Gambey}, a description of which is to follow in the next volume.
